\documentclass[11pt]{article}

\usepackage[margin=1in]{geometry}
\usepackage{amsmath,amssymb,amsthm}

\newtheorem{theorem}{Theorem}[section]
\newtheorem{lemma}[theorem]{Lemma}
\newtheorem{proposition}[theorem]{Proposition}
\newtheorem{corollary}[theorem]{Corollary}
\theoremstyle{definition}
\newtheorem{definition}[theorem]{Definition}
\theoremstyle{remark}
\newtheorem{remark}[theorem]{Remark}

\newcommand{\area}{\operatorname{area}}
\newcommand{\dinv}{\operatorname{dinv}}
\newcommand{\defc}{\operatorname{defc}}
\newcommand{\Cat}{\operatorname{Cat}}

\title{The Modulo-One Partition Formula}
\author{}
\date{}

\begin{document}
\maketitle

\section{Setup}

Fix
\[
  \tau\ge2,\qquad s\ge6,\qquad r=\tau s+1,
  \qquad M=\tau\binom{s}{2},
\]
and put
\[
  B=\tau(s-2)-1.
\]
A normalized step-$\tau$ Dyck word is a word
\[
  D=(x_0,\ldots,x_{s-1})
\]
of nonnegative integers satisfying
\[
  x_0=0,
  \qquad x_{i+1}\le x_i+\tau.
\]
Its statistics are
\[
 \area(D)=\sum_i x_i,
 \qquad
 \dinv_\tau(D)=\sum_{i<j}d_\tau(x_i,x_j),
 \qquad
 \defc_\tau(D)=M-\area(D)-\dinv_\tau(D),
\]
where
\[
 d_\tau(u,v)=
 \begin{cases}
   (u+\tau-v)_+,&u\le v,\\
   (v+\tau+1-u)_+,&u>v.
 \end{cases}
\]

An occurrence $x_j=e>0$ is extractable if exactly one earlier occurrence
lies in
\[
  [\max(0,e-\tau),e)
\]
and its deletion leaves a normalized word.  A \emph{full skeleton} is a
normalized word with no extractable occurrence.

We use three root sets:
\begin{align*}
 \mathcal F_B
   &=\{S:S\text{ is a full skeleton and }\defc_\tau(S)\le B\},\\
 \mathcal Z_B
   &=\{Z:Z\text{ is normalized},\ Z_0=Z_1=0,
                    \ \defc_\tau(Z)\le B\},\\
 \mathcal P_B^{(s-2)}
   &=\{\lambda:\lambda\text{ is a partition},\
          \lambda_1\le s-2,\ |\lambda|\le B\}.
\end{align*}
The empty partition is included in the last set.

We shall use the modulo-one string-decomposition theorem proved in
\texttt{mod1\_decomp.tex}.  Thus we assume the $q,t$ symmetry of
$\Cat_{\tau s+1,s}(q,t)$ and the East5 and West5 nonfailure hypotheses made
there.  We will prove here the remaining special-skeleton area bound, rather
than assume it.

\section{Full skeletons and initial double zeros}

Put
\[
 K_\tau(u,v)=\tau-d_\tau(u,v)
 =
 \begin{cases}
   \min(\tau,v-u),&u\le v,\\
   \min(\tau,u-v-1),&u>v.
 \end{cases}
\]
Since there are $\binom{s}{2}$ pairs, we have the useful identity
\begin{equation}
 \defc_\tau(D)
 =\sum_{i<j}K_\tau(x_i,x_j)-\sum_jx_j.
 \label{eq:defect-kernel}
\end{equation}

\begin{lemma}[Large-entry bound]
\label{lem:large-entry}
Let $D=(0,0,x_2,\ldots,x_{s-1})$ be normalized.  If some entry of $D$ is
greater than $\tau$, then
\[
  \defc_\tau(D)\ge\tau(s-2).
\]
\end{lemma}

\begin{proof}
For $1\le j\le s-1$, form a column whose first $x_j$ cells are $1$s and
whose next $\tau$ cells are $0$s.  A $0$ and a $1$ are called a contact if
either the $0$ is strictly to the left of the $1$ in the same row, or the
$1$ is one row below the $0$ and strictly to its left.  For $i<j$, the
number of contacts using columns $i$ and $j$ is
\[
 \begin{cases}
  \min(\tau,x_j-x_i),&x_i\le x_j,\\
  \min(\tau,x_i-x_j-1),&x_i>x_j,
 \end{cases}
\]
which is $K_\tau(x_i,x_j)$.  If $C(D)$ is the total number of contacts and
\[
  u(D)=\sum_{j=1}^{s-1}(x_j-\tau)_+,
\]
then the pairs involving $x_0=0$ in \eqref{eq:defect-kernel} give
\begin{equation}
  \defc_\tau(D)=C(D)-u(D).
  \label{eq:contacts}
\end{equation}

Choose $k$ with $x_k>\tau$.  Consider the $\tau(s-2)$ cells in the first
$\tau$ rows and columns $2,\ldots,s-1$.  A cell containing $1$ is in contact
with the $0$ in column $1$ of the same row.  If the cell is a $0$ in a
column left of $k$, it is in contact with the $1$ in column $k$ of the same
row.  If it is a $0$ in a column right of $k$, it is in contact with the $1$
in column $k$ one row below.  These are $\tau(s-2)$ distinct contacts.

There is also a distinct contact for each of the $u(D)$ ones below row
$\tau$.  For such a $1$ in row $a$ and column $j$, choose $i<j$ maximal
with $x_i<a$; this exists because $x_1=0$.  Then $x_{i+1}\ge a$, and
normalization gives
\[
  x_i<a\le x_{i+1}\le x_i+\tau.
\]
Thus column $i$ contains a $0$ in row $a$, and it is in contact with the
chosen $1$ in the same row.  These contacts lie below row $\tau$, so none
was counted previously.  Hence
\[
  C(D)\ge\tau(s-2)+u(D).
\]
Equation~\eqref{eq:contacts} proves the result.
\end{proof}

\begin{proposition}[The two descriptions of the roots]
\label{prop:full-zero}
In the range $\defc_\tau\le B$,
\[
  D\text{ is a full skeleton}
  \quad\Longleftrightarrow\quad
  D\text{ begins with }(0,0).
\]
Consequently $\mathcal F_B=\mathcal Z_B$.
\end{proposition}

\begin{proof}
First suppose that $D$ is full.  If $x_1>0$, then $x_1\le\tau$ and $x_0=0$
is its unique predecessor.  Since $x_1$ is not extractable, its outgoing
splice must fail, so $x_2>x_0+\tau$.  Normalization then makes $x_2$ a new
record with $x_1$ as its unique predecessor.  Repeating the same argument
produces a strictly increasing record chain.  The final member of this chain
has a valid outgoing splice, or lies at the right boundary, and is therefore
extractable, a contradiction.  Hence $x_1=0$.

Conversely, suppose $D$ begins with $(0,0)$ and
$\defc_\tau(D)\le B<\tau(s-2)$.  Lemma~\ref{lem:large-entry} shows that every
entry is at most $\tau$.  If $e>0$ occurs, both initial zeros lie in
\[
  [\max(0,e-\tau),e)=[0,e).
\]
Thus $e$ has at least two predecessors and is not extractable.  Zero itself
is never extractable, so $D$ is full.
\end{proof}

\section{Bounded partitions}

We now use the one external bijection needed for the partition form of the
answer.

\begin{theorem}[Hawkes's bounded-partition bijection]
\label{thm:partition-bijection}
There is a bijection
\[
  \Phi:\mathcal P_B^{(s-2)}\longrightarrow\mathcal Z_B
\]
such that
\begin{equation}
  \defc_\tau(\Phi(\lambda))=|\lambda|,
  \qquad
  \area(\Phi(\lambda))=\ell(\lambda).
  \label{eq:partition-statistics}
\end{equation}
\end{theorem}

\begin{proof}[Source]
This is the composition of the partition--matrix bijection, the matrix
involution interchanging height and width, and the matrix--maximal-path
bijection in Section~5 of~\cite{Hawkes2024}.  In that paper set
$\ell=s-1$ and $m=\tau$.  Its maximal paths are precisely the normalized
position-coordinate words beginning with $(0,0)$, and its condition
$|\lambda|<(\ell-1)m$ is exactly $|\lambda|\le B$.
\end{proof}

\begin{corollary}[Skeleton position]
\label{cor:skeleton-position}
Every full skeleton $S$ with $\defc_\tau(S)\le B$ satisfies
\[
  \area(S)\le\defc_\tau(S).
\]
In particular it lies strictly below the middle of its defect layer.
\end{corollary}

\begin{proof}
By Proposition~\ref{prop:full-zero} and
Theorem~\ref{thm:partition-bijection}, write $S=\Phi(\lambda)$.  Since all
parts of a partition are positive,
\[
  \area(S)=\ell(\lambda)\le|\lambda|=\defc_\tau(S).
\]
The strict lower-half assertion follows from the calculation in
\texttt{mod1\_decomp.tex}; explicitly, if $d=\defc_\tau(S)\le B$, then
\[
  2\area(S)+d\le3d\le3B<M.
\]
\end{proof}

\section{The three forms of the decomposition}

\begin{theorem}[Full-skeleton, double-zero, and partition forms]
\label{thm:three-forms}
Assume the $q,t$ symmetry and the East5 and West5 nonfailure hypotheses used
in \texttt{mod1\_decomp.tex}.  The weak lower half of every defect layer
$d\le B$ is partitioned into strings whose roots may be indexed equivalently
by the defect-$d$ subsets of
\[
  \mathcal F_B,
  \qquad \mathcal Z_B,
  \qquad \mathcal P_B^{(s-2)}.
\]
For the partition set, its defect-$d$ subset means the partitions of size
$d$.
For $\lambda\in\mathcal P_B^{(s-2)}$, the root is
$S_\lambda=\Phi(\lambda)$, and its lower-half string has one word at every
area
\[
  \ell(\lambda),\ \ell(\lambda)+1,\ldots,
  \left\lfloor\frac{M-|\lambda|}{2}\right\rfloor.
\]

Consequently the low-defect part of the rational $q,t$-Catalan polynomial
has the following three equal root-indexed forms:
\begin{align}
 &\sum_{\substack{D\text{ normalized}\\\defc_\tau(D)\le B}}
 q^{\dinv_\tau(D)}t^{\area(D)}
 \notag\\[1mm]
 &=\sum_{S\in\mathcal F_B}
 \frac{
 q^{\dinv_\tau(S)+1}t^{\area(S)}
 -q^{\area(S)}t^{\dinv_\tau(S)+1}
 }{q-t}
 \label{eq:full-form}\\[1mm]
 &=\sum_{Z\in\mathcal Z_B}
 \frac{
 q^{\dinv_\tau(Z)+1}t^{\area(Z)}
 -q^{\area(Z)}t^{\dinv_\tau(Z)+1}
 }{q-t}
 \label{eq:zero-form}\\[1mm]
 &=\sum_{\lambda\in\mathcal P_B^{(s-2)}}
 \frac{
 q^{M-|\lambda|-\ell(\lambda)+1}t^{\ell(\lambda)}
 -q^{\ell(\lambda)}t^{M-|\lambda|-\ell(\lambda)+1}
 }{q-t}.
 \label{eq:partition-form}
\end{align}
\end{theorem}

\begin{proof}
The exceptional full skeleton from \texttt{mod1\_decomp.tex} has defect
$\tau(s-2)=B+1$.  Hence in the present range the special skeletons appearing
there are exactly the full skeletons in $\mathcal F_B$.
Corollary~\ref{cor:skeleton-position} proves the skeleton-position hypothesis
of that decomposition theorem.  Applying it gives the lower-half strings and
\eqref{eq:full-form}.

Proposition~\ref{prop:full-zero} identifies $\mathcal F_B$ with
$\mathcal Z_B$, giving \eqref{eq:zero-form}.  Finally reindex by the bijection
$\Phi$.  Equation~\eqref{eq:partition-statistics} and the identity
\[
 \dinv_\tau(S_\lambda)
 =M-\defc_\tau(S_\lambda)-\area(S_\lambda)
 =M-|\lambda|-\ell(\lambda)
\]
give both the stated lower-half area interval and
\eqref{eq:partition-form}.
\end{proof}

\begin{thebibliography}{9}

\bibitem{Hawkes2024}
Graham Hawkes.
\newblock \emph{A conjectured formula for the rational $q,t$-Catalan
polynomial}.
\newblock Annals of Combinatorics 28 (2024), 749--795;
arXiv:2208.00577.

\end{thebibliography}

\end{document}
